% Source-derived chapter 04: Conjecture~\ref{conjnum} for hyper-K\"ahler fourfolds
% Original source: riemann-roch-polynomials-hyperkahler-fourfolds
\chapter{Conjecture~\ref{conjnum} for hyper-K\"ahler fourfolds}
\label{riemann-roch-polynomials-hyperkahler-fourfolds-chapter-04}
\source{riemann-roch-polynomials-hyperkahler-fourfolds}

\begin{remark}[Source section]
\label{riemann-roch-polynomials-hyperkahler-fourfolds-chapter-04-source}
\source{riemann-roch-polynomials-hyperkahler-fourfolds}
\source{riemann-roch-polynomials-hyperkahler-fourfolds}
This chapter reproduces the corresponding section of the retrieved
LaTeX source, including its definitions, statements, examples, and
proofs. It has not yet been translated into Lean.
\notready
\end{remark}

% The source section title was: Conjecture~\ref{conjnum} for hyper-K\"ahler fourfolds
\label{sec:IsotropicClasses}
\source{riemann-roch-polynomials-hyperkahler-fourfolds}

The main result of  this section is the proof of  Conjecture~\ref{conjnum} in dimension $4$ (see Theorem~\ref{th43a}).\ It is a simple consequence of the work \cite{gua} of Guan  who gave a list of possible Betti numbers for hyper-K\"ahler fourfolds.

Let $X$ be a \hk\ fourfold.\ Following \cite[Section~2.4]{jia}, we
 set
\begin{equation}\label{ax}
\source{riemann-roch-polynomials-hyperkahler-fourfolds}
A_X\coloneqq\int_X\td^{1/2}(X)= \frac 1 {5760} ( 7c_2^2(X) -4c_4(X) ) = \frac 18 \Bigl(7-\frac 1 {432} c_4(X)\Bigr).
\end{equation}
It is known that $c_4(X)$ is divisible by 12 (see, for example, \cite[Proposition 2.4]{gr}), hence $288 A_X$ is an integer.

\begin{lemm}\label{lemmguan}
\source{riemann-roch-polynomials-hyperkahler-fourfolds}
Let $X$ be a \hk\ fourfold.\ Then,
\begin{enumerate}[{\rm (a)}]
\item\label{aax} either $b_2(X) =23$,  $b_3(X) =0$, and the Hodge numbers of $X$ are those of the Hilbert square of a K3 surface, in which case $ c_4(X)= 324$ and $A_X=\frac{25}{32}$,
\source{riemann-roch-polynomials-hyperkahler-fourfolds}
\item\label{bax} or $b_2(X) \le 8$, in which case $c_4(X)\le 144$ and  $\frac{5}{6}\le A_X\le \frac{131}{144}$.
\source{riemann-roch-polynomials-hyperkahler-fourfolds}
 \end{enumerate}
In particular, if $t\in [0,\frac13)$, then $4A_X-t $ is an integer only when $t=\frac18$ and $A_X=\frac{25}{32}$.
\end{lemm}

\begin{proof}
This follows from~\eqref{ax}, the relation $c_4(X)=3(4b_2(X)+16-b_3(X))$, and the list of possible Betti numbers given in  \cite[Main Theorem]{gua}.
\end{proof}


We now introduce classes $\lll,\mm\in H^2(X,\Z)$ such that $q_X(\lll)=0$ and define $a=\frac12\int_X\lll^2\mm^2$ as in~\eqref{defa2}.\ By Lemma~\ref{lemmaa} and~\eqref{defia}, it is a  nonnegative integer.



\begin{lemm}\label{lemmax}
\source{riemann-roch-polynomials-hyperkahler-fourfolds}
The number $\sqrt{2aA_X} $ is   rational.
\end{lemm}

\begin{proof}
By~\eqref{defia}, we have $c_Xq_X(\lll,\mm)^2=3a$.\ Using Section~\ref{subsec:RRpolynomial} and~\cite[Lemma~5.7]{jia}, we obtain
\begin{equation}\label{prr}
\source{riemann-roch-polynomials-hyperkahler-fourfolds}
   P_{RR,X}(T)=\frac{c_X}{24}T^2+ T\sqrt{\frac23c_XA_X} +3=\frac{a}{8}\Bigl(\frac{T}{q_X(\lll,\mm)}\Bigr)^2+ \sqrt{2aA_X}\Bigl(\frac{T}{q_X(\lll,\mm)}\Bigr) +3.
\end{equation}
This implies that $\sqrt{2aA_X} $ is a rational number.
\end{proof}

We are now ready to prove (a stronger form of) Conjecture~\ref{conjnum} (which corresponds to the case $a=1$) in dimension 4.\ Analogous, but weaker,  results for low values $a$ will be given in Theorem~\ref{th43}.

\begin{theo}\label{th43a}
\source{riemann-roch-polynomials-hyperkahler-fourfolds}
Let $X$ be a  \hk\ fourfold with classes~$\lll,\mm\in H^2(X,\Z)$ such that $\int_X \lll^{4}=0$
and   $\int_X \lll^{2}\mm^{2}=2$.\
 The Chern and Hodge numbers of $X$ are those of the Hilbert square of a K3 surface  and all the  conclusions of Theorem~\ref{th21} hold.
\end{theo}

\begin{proof}
Changing   $\mm$ into $\pm\mm+r\lll$, for some $r\in\Z$, if necessary, we may assume $-q_X(\lll,\mm)< q_X(\mm)\le q_X(\lll,\mm)$ or, equivalently, $\gamma\coloneqq \frac{q_X(\mm)}{q_X(\lll,\mm)}\in (-1,1] $.\

As in~\eqref{defP}, we introduce the polynomial
$$P(k)\coloneqq P_{RR,X}(q_X(k\lll+ \mm))=P_{RR,X}(2kq_X(\lll,\mm)+q_X(\mm)).$$
Using~\eqref{prr}, we compute
\begin{align*}
 P(k)&= \frac{1}{8}(2k+\gamma)^2+ \sqrt{2A_X}(2k+\gamma) +3  \\
 &=\frac{1}{2}k^2+\Bigl( \frac{1}{2}\gamma+2\sqrt{2A_X} \Bigr)k+\frac{1}{8}\gamma^2+ \gamma\sqrt{2A_X}+3\\
 &=:\frac{1}{2}k^2 +  b  k +c.
\end{align*}
Since $P$ takes integral values on integers, $\frac{1}2+b=P(1)-P(0)$ and $c=P(0)$ are integers; we write $b=\frac12+b'$, with $b'\in\Z$.

We also note that
\begin{align*}
    4A_X-\frac{b^2}{2}&= 4A_X-\frac{1}{2}\Bigl( \frac{1}{2}\gamma+2\sqrt{2A_X} \Bigr)^2\nonumber\\
    &=4A_X-\frac{1}{2}\Bigl( \frac{1}{4}\gamma^2+2\gamma\sqrt{2A_X}+8A_X \Bigr)\\
    &=3-c\nonumber
\end{align*}
is an integer, hence so is $4A_X-\frac{1}{8}$.\ By Lemma~\ref{lemmguan}, this is only possible when $A_X=\frac{25}{32}$ and the Chern and Hodge numbers of $X$ are those of the Hilbert square of a K3 surface.\ We also have
$$\frac12+b'=b=\frac{1}{2}\gamma+\frac52,$$
which implies that $\gamma$ is an even integer.\ Since $\gamma \in (-1,1] $, we obtain $\gamma=0$,  hence $q_X(\mm)=0$, $b=\frac52$, and $c=3$.\ By  the very definition~\eqref{defP} of $P$, one has
\[
P_{RR,X}(2kq_X(\lll,\mm))=P(k)= \frac{1}{2}(k^2 +  5  k +6)=\binom{k+3}{2}.
\]
By Lemma~\ref{star} (applied with $c=2$, $c'=1$, and $q=2q_X(\lll,\mm)$), we get $q_X(\lll,\mm)=1$, the quadratic form $q_X$ is even, and all the conclusions of Theorem~\ref{th21} hold.
\end{proof}


%%%%%%%%%%%%%%%%%%%%%
